% GENERATED FILE. Edit canonical lessons, then run npm run generate:books.
% canonical-source-hash: fnv1a64:00387faeb9dd658d
% canonical-lessons: SA-C167-kacchapah, SA-C167-mahisi, SA-C167-vatsah, SA-C167-aja, SA-C167-asva, SA-R167-first-pass-words-from-chapters-159-162, SA-R167-second-pass-words-from-chapters-163-167

\chapter{Turtle, She-buffalo, Calf, She-goat, Mare}
\label{ch:sa-turtle-she-buffalo-calf-she-goat-mare}
\hlchaptermodality{167}

\begin{chapteropening}
\textbf{By the end of this chapter:} \emph{I can say a turtle, a she-buffalo, a calf, a she-goat and a mare.}

Complete the last lesson of chapter 167: I can say a turtle, a she-buffalo, a calf, a she-goat and a mare.
\end{chapteropening}

\section[\texorpdfstring{\sk{कच्छपः} (kacchapaḥ)}{कच्छपः (kacchapaḥ)}]{\sk{कच्छपः} (kacchapaḥ) — a turtle}

\label{lesson:SA-C167-kacchapah}



\begin{quote}
Before the new one: say the Sanskrit for a banyan tree, then the Sanskrit for a palm tree.
\end{quote}

\subsection*{You'll want to know: \sk{कच्छपः}}
\textbf{\sk{कच्छपः}} — \emph{kacchapaḥ} — \textquotedblleft{}a turtle\textquotedblright{}.

\begin{grammarlens}[title={What you've built}]
One more word for this chapter.
\end{grammarlens}

\subsection*{Your turn}
\begin{itemize}
\raggedright
  \item \emph{Say it:} \emph{kacchapaḥ}
  \item \emph{Say it:} \emph{kacchapaḥ}, once more
  \item \emph{Say it:} say \emph{kacchapaḥ}
  \item \emph{Recall:} say the Sanskrit for a banyan tree, then the Sanskrit for a palm tree, then say \emph{kacchapaḥ} again
\end{itemize}

\begin{tcolorbox}[breakable,colback=teal!4,colframe=teal!35!black,title={Before you move on}]
What does \sk{कच्छपः} mean? (\textquotedblleft{}A turtle\textquotedblright{}.) Say it once more.
\end{tcolorbox}
\section[\texorpdfstring{\sk{महिषी} (mahiṣī)}{महिषी (mahiṣī)}]{\sk{महिषी} (mahiṣī) — a she-buffalo}

\label{lesson:SA-C167-mahisi}



\begin{quote}
Before the new one: say the Sanskrit for a palm tree, then the Sanskrit for a turtle.
\end{quote}

\subsection*{You'll want to know: \sk{महिषी}}
\textbf{\sk{महिषी}} — \emph{mahiṣī} — \textquotedblleft{}a she-buffalo\textquotedblright{}.

\begin{grammarlens}[title={What you've built}]
One more word for this chapter.
\end{grammarlens}

\subsection*{Your turn}
\begin{itemize}
\raggedright
  \item \emph{Say it:} \emph{mahiṣī}
  \item \emph{Say it:} \emph{mahiṣī}, once more
  \item \emph{Say it:} say \emph{mahiṣī}
  \item \emph{Recall:} say the Sanskrit for a palm tree, then the Sanskrit for a turtle, then say \emph{mahiṣī} again
\end{itemize}

\begin{tcolorbox}[breakable,colback=teal!4,colframe=teal!35!black,title={Before you move on}]
What does \sk{महिषी} mean? (\textquotedblleft{}A she-buffalo\textquotedblright{}.) Say it once more.
\end{tcolorbox}
\section[\texorpdfstring{\sk{वत्सः} (vatsaḥ)}{वत्सः (vatsaḥ)}]{\sk{वत्सः} (vatsaḥ) — a calf}

\label{lesson:SA-C167-vatsah}



\begin{quote}
Before the new one: say the Sanskrit for a turtle, then the Sanskrit for a she-buffalo.
\end{quote}

\subsection*{You'll want to know: \sk{वत्सः}}
\textbf{\sk{वत्सः}} — \emph{vatsaḥ} — \textquotedblleft{}a calf\textquotedblright{}.

\begin{grammarlens}[title={What you've built}]
One more word for this chapter.
\end{grammarlens}

\subsection*{Your turn}
\begin{itemize}
\raggedright
  \item \emph{Say it:} \emph{vatsaḥ}
  \item \emph{Say it:} \emph{vatsaḥ}, once more
  \item \emph{Say it:} say \emph{vatsaḥ}
  \item \emph{Recall:} say the Sanskrit for a turtle, then the Sanskrit for a she-buffalo, then say \emph{vatsaḥ} again
\end{itemize}

\begin{tcolorbox}[breakable,colback=teal!4,colframe=teal!35!black,title={Before you move on}]
What does \sk{वत्सः} mean? (\textquotedblleft{}A calf\textquotedblright{}.) Say it once more.
\end{tcolorbox}
\section[\texorpdfstring{\sk{अजा} (ajā)}{अजा (ajā)}]{\sk{अजा} (ajā) — a she-goat}

\label{lesson:SA-C167-aja}



\begin{quote}
Before the new one: say the Sanskrit for a she-buffalo, then the Sanskrit for a calf.
\end{quote}

\subsection*{You'll want to know: \sk{अजा}}
\textbf{\sk{अजा}} — \emph{ajā} — \textquotedblleft{}a she-goat\textquotedblright{}.

\begin{grammarlens}[title={What you've built}]
One more word for this chapter.
\end{grammarlens}

\subsection*{Your turn}
\begin{itemize}
\raggedright
  \item \emph{Say it:} \emph{ajā}
  \item \emph{Say it:} \emph{ajā}, once more
  \item \emph{Say it:} say \emph{ajā}
  \item \emph{Recall:} say the Sanskrit for a she-buffalo, then the Sanskrit for a calf, then say \emph{ajā} again
\end{itemize}

\begin{tcolorbox}[breakable,colback=teal!4,colframe=teal!35!black,title={Before you move on}]
What does \sk{अजा} mean? (\textquotedblleft{}A she-goat\textquotedblright{}.) Say it once more.
\end{tcolorbox}
\section[\texorpdfstring{\sk{अश्वा} (aśvā)}{अश्वा (aśvā)}]{\sk{अश्वा} (aśvā) — a mare}

\label{lesson:SA-C167-asva}



\begin{quote}
Before the new one: say the Sanskrit for a calf, then the Sanskrit for a she-goat.
\end{quote}

\subsection*{You'll want to know: \sk{अश्वा}}
\textbf{\sk{अश्वा}} — \emph{aśvā} — \textquotedblleft{}a mare\textquotedblright{}.

\begin{grammarlens}[title={What you've built}]
One more word for this chapter.
\end{grammarlens}

\subsection*{Your turn}
\begin{itemize}
\raggedright
  \item \emph{Say it:} \emph{aśvā}
  \item \emph{Say it:} \emph{aśvā}, once more
  \item \emph{Say it:} say \emph{aśvā}
  \item \emph{Recall:} say the Sanskrit for a calf, then the Sanskrit for a she-goat, then say \emph{aśvā} again
\end{itemize}

\begin{tcolorbox}[breakable,colback=teal!4,colframe=teal!35!black,title={Before you move on}]
What does \sk{अश्वा} mean? (\textquotedblleft{}A mare\textquotedblright{}.) Say it once more.
\end{tcolorbox}
\section[(dialogue)]{Words from chapters 159-162 — a first pass}

\label{lesson:SA-R167-first-pass-words-from-chapters-159-162}



\begin{quote}
Say the words for a she-goat and a mare.
\end{quote}

\subsection*{Your turn}
Say these aloud:

\begin{itemize}
\raggedright
  \item cumin, mustard, asafoetida, sesame, a chickpea
  \item a meal, a curry, food, eating, a small pitcher
  \item a conch, a flute, a mridanga drum, an anklet, an earring
  \item wages, a debt, profit, loss, trade
\end{itemize}

\begin{tcolorbox}[breakable,colback=teal!4,colframe=teal!35!black,title={Before you move on}]
Cumin? (\textbf{\sk{जीरकम्}}, \emph{jīrakam}.) Mustard? (\textbf{\sk{सर्षपः}}, \emph{sarṣapaḥ}.) Asafoetida? (\textbf{\sk{हिंगु}}, \emph{hiṅgu}.) Sesame? (\textbf{\sk{तिलः}}, \emph{tilaḥ}.) A chickpea? (\textbf{\sk{चणकः}}, \emph{caṇakaḥ}.) A meal? (\textbf{\sk{भोजनम्}}, \emph{bhojanam}.) A curry? (\textbf{\sk{व्यञ्जनम्}}, \emph{vyañjanam}.) Food? (\textbf{\sk{आहारः}}, \emph{āhāraḥ}.) Eating? (\textbf{\sk{अशनम्}}, \emph{aśanam}.) A small pitcher? (\textbf{\sk{कलशी}}, \emph{kalaśī}.) A conch? (\textbf{\sk{शंखः}}, \emph{śaṅkhaḥ}.) A flute? (\textbf{\sk{वेणुः}}, \emph{veṇuḥ}.) A mridanga drum? (\textbf{\sk{मृदंगः}}, \emph{mṛdaṅgaḥ}.) An anklet? (\textbf{\sk{नूपुरम्}}, \emph{nūpuram}.) An earring? (\textbf{\sk{कुण्डलम्}}, \emph{kuṇḍalam}.) Wages? (\textbf{\sk{वेतनम्}}, \emph{vetanam}.) A debt? (\textbf{\sk{ऋणम्}}, \emph{ṛṇam}.) Profit? (\textbf{\sk{लाभः}}, \emph{lābhaḥ}.) Loss? (\textbf{\sk{हानिः}}, \emph{hāniḥ}.) Trade? (\textbf{\sk{व्यापारः}}, \emph{vyāpāraḥ}.)
\end{tcolorbox}
\section[(dialogue)]{Words from chapters 163-167 — a second pass}

\label{lesson:SA-R167-second-pass-words-from-chapters-163-167}



\begin{quote}
Say the words for sport and gambling.
\end{quote}

\subsection*{Your turn}
Say these aloud:

\begin{itemize}
\raggedright
  \item sport, gambling, a festival, a wedding, a gift
  \item a courtyard, a small garden, a park, a playground, an airport
  \item dust, smoke, a flame, a live coal, iron
  \item a sacred fig tree, a neem tree, a mango tree, a banyan tree, a palm tree
  \item a turtle, a she-buffalo, a calf, a she-goat, a mare
\end{itemize}

\begin{tcolorbox}[breakable,colback=teal!4,colframe=teal!35!black,title={Before you move on}]
Sport? (\textbf{\sk{क्रीडा}}, \emph{krīḍā}.) Gambling? (\textbf{\sk{द्यूतम्}}, \emph{dyūtam}.) A festival? (\textbf{\sk{उत्सवः}}, \emph{utsavaḥ}.) A wedding? (\textbf{\sk{विवाहः}}, \emph{vivāhaḥ}.) A gift? (\textbf{\sk{उपहारः}}, \emph{upahāraḥ}.) A courtyard? (\textbf{\sk{प्रांगणम्}}, \emph{prāṅgaṇam}.) A small garden? (\textbf{\sk{वाटिका}}, \emph{vāṭikā}.) A park? (\textbf{\sk{उपवनम्}}, \emph{upavanam}.) A playground? (\textbf{\sk{क्रीडांगणम्}}, \emph{krīḍāṅgaṇam}.) An airport? (\textbf{\sk{विमानपत्तनम्}}, \emph{vimānapattanam}.) Dust? (\textbf{\sk{धूलिः}}, \emph{dhūliḥ}.) Smoke? (\textbf{\sk{धूमः}}, \emph{dhūmaḥ}.) A flame? (\textbf{\sk{ज्वाला}}, \emph{jvālā}.) A live coal? (\textbf{\sk{अंगारः}}, \emph{aṅgāraḥ}.) Iron? (\textbf{\sk{लौहम्}}, \emph{lauham}.) A sacred fig tree? (\textbf{\sk{अश्वत्थः}}, \emph{aśvatthaḥ}.) A neem tree? (\textbf{\sk{निम्बः}}, \emph{nimbaḥ}.) A mango tree? (\textbf{\sk{आम्रवृक्षः}}, \emph{āmravṛkṣaḥ}.) A banyan tree? (\textbf{\sk{वटः}}, \emph{vaṭaḥ}.) A palm tree? (\textbf{\sk{तालः}}, \emph{tālaḥ}.) A turtle? (\textbf{\sk{कच्छपः}}, \emph{kacchapaḥ}.) A she-buffalo? (\textbf{\sk{महिषी}}, \emph{mahiṣī}.) A calf? (\textbf{\sk{वत्सः}}, \emph{vatsaḥ}.) A she-goat? (\textbf{\sk{अजा}}, \emph{ajā}.) A mare? (\textbf{\sk{अश्वा}}, \emph{aśvā}.)
\end{tcolorbox}
